% Candidate bridge for the full BFJ psh class, not a change of definition.
% Requires independent and production Rethlas review before integration.
\begin{lemma}
\label{lem:ch23-psh-vertex-chord}
Assume the standing BFJ hypotheses of this chapter and \(\dim X=1\).
Let \(\mathcal X\) be a strictly semistable model determining
\(\theta\), write \(\Gamma=\Delta_{\mathcal X}\), and let \(V\)
be its vertex set.  For \(\varphi\in\PSH(X,\theta)\), let
\(A\varphi\) be the function that agrees with \(\varphi\) on
\(V\), is affine on every edge of \(\Gamma\), and is extended to
\(X^{\An}\) using the retraction \(p_{\mathcal X}\).
Then \(A\varphi\) is a real model \(\theta\)-psh function determined
on \(\mathcal X\), and
\[
                 \varphi\leq A\varphi.
\]
If a real model obstacle \(u\) is determined on \(\mathcal X\) and
\(\varphi\leq u\), then \(A\varphi\leq u\) as well.
No rational-model assertion is made about \(A\varphi\).
\end{lemma}

\begin{proof}
If the strictly semistable model has a node of thickness greater than
one, first take its regular semistable resolution over the original
discretely valued field.  Locally this resolves \(xy=\pi^a\) by
subdividing its length-\(a\) edge into unit-length edges; the construction
is compatible with etale node charts.  The metric skeleton and its
retraction are unchanged, while extra vertices are inserted.
The pullback of \(\theta\) is determined on this regular SNC model.
All uses below of the BFJ fixed-skeleton approximation and retraction
inequality are made on that fixed SNC determination, then restricted
to the same metric graph with its original vertex set.  No refinement
is chosen after an approximant has been selected.  The chord still
uses the original vertices; its real Cartier interpretation on the
original model is checked below.

We first check the assertion for a \(\theta\)-psh real model
function \(\chi\).  The edge lengths are the analytic modulus
lengths and the edge weights are the residue degrees of the nodes.
We use these weights at both ends of every edge, as in
\eqref{eq:ch31-graph-laplacian}.

Pass temporarily to the completed algebraic closure, and work
component by component if the curve becomes disconnected.  A common
strictly semistable refinement determines the finitely many real
model classes defining \(\theta\) and \(\chi\); the semistable
domination input is \cite[\S7]{BL85}, as used in
\cite[proof of Proposition~3.7]{GJKM19}.  On its skeleton \(\Gamma'\),
the curvature of \(\theta+\mathrm d\mathrm d^c\chi\) is
\[
                      D'+\Delta\chi\geq0.
\]
This follows from numerical nefness and the component-degree formula,
extended real-linearly; it does not require \(\chi\) to be rational.
The pullback of \(\theta\) has degree zero on the components
contracted by the refinement, so \(D'\) is supported on the old
vertices.  If \(r:\Gamma'\to\Gamma\) denotes retraction to the
base-changed old skeleton, then
\[
      D+\Delta(\chi|_\Gamma)
          =r_*(D'+\Delta\chi)\geq0.
\]
This is pushforward from a refinement of graphs over the same field,
not an inference of positivity from a scalar pushforward in the
opposite direction.  The old skeleton may first be subdivided;
the identity is unchanged by subdivision.
Here is the graph verification of the identity for this possibly
nonfactoring \(\chi\).  By \cite[Proposition~2.2.26(iii)]{Thu05},
each new connected piece attaches to the old skeleton at one point.
Subtract \((\chi|_\Gamma)\circ r\) from \(\chi\).  The difference
vanishes on \(\Gamma\).  On each attached piece its Laplacian,
including the mass at the attaching point, has total mass zero by
cancellation of outgoing slopes along its edges.  Retraction sends
all these masses to the attaching point, so their pushforward is zero.
For the subtracted pullback function, the pushforward identity is
\cite[Proposition~2.2.26(iv)]{Thu05}.  Adding the two identities proves
the displayed formula without assuming that \(\chi\) factors
through the original retraction.

In particular \(\chi\) is convex on each old open edge.  Replacing
it there by the chord of its endpoint values increases both outgoing
endpoint slopes.  The chord has no interior slope jump.  Its
curvature is therefore nonnegative at the old vertices and zero
elsewhere on the old open edges.  These are exactly the numerical
nef inequalities for the real model class of \(A\chi\), again by
\cite[Proposition~3.7]{GJKM19}.  The chord is the pullback of the
original-field real model function with the specified vertex values.
Nefness of a real model numerical class reflects under scalar
extension, by the degree and projection-formula argument in
\cref{thm:ch29-real-class-curve-model-envelope}.  It follows that
\(A\chi\) is \(\theta\)-psh over the original field.

For completeness, arbitrary real vertex values do define a real
model function on this strictly semistable model.  On a regular
model this is the vertical real Cartier divisor whose coefficients
are those values.  At a node with local equation \(xy=\pi^a\),
each component is \(\mathbf Q\)-Cartier: a multiple \(a\) of
one branch divisor is the divisor of \(x\) or \(y\).
The same assertion is checked on an etale node chart when the node
is not rational over the residue field.  Thus a real linear
combination of components is real Cartier.  Its model function
has precisely the edge-affine interpolation of the vertex values.
Equivalently one may compute it on the regular subdivision of that
node, where the new coefficients are the interpolated values.

Now let \(\varphi\) be an arbitrary BFJ-psh function.  By its
defining approximation on the fixed skeleton, there are
\(\theta\)-psh real model functions \(\chi_j\) with
\(\chi_j|_\Gamma\to\varphi|_\Gamma\) uniformly.  Only this
fixed-skeleton approximation is used.  The vertex vectors of
\(A\chi_j\) converge to those of \(A\varphi\).
There are finitely many numerical inequalities at the vertices;
they are affine and hence closed in these vertex values.  Since
they hold for every \(A\chi_j\), they hold for \(A\varphi\).
Thus \(A\varphi\) is a real model \(\theta\)-psh function.

Finally the BFJ retraction inequality gives
\(\varphi\leq\varphi\circ p_{\mathcal X}\), and edgewise
convexity gives \(\varphi|_\Gamma\leq(A\varphi)|_\Gamma\).
This proves \(\varphi\leq A\varphi\) globally.  If
\(\varphi\leq u\), the endpoint values of \(A\varphi\) are
bounded above by those of the edge-affine function \(u\).
Affine interpolation and retraction then give \(A\varphi\leq u\)
globally, as asserted.
\end{proof}
